\documentclass[10pt,a4paper]{amsart}
\usepackage[margin=19mm]{geometry}
\usepackage{amsmath,amssymb,mathtools}
\usepackage[scr=rsfs,cal=euler]{mathalfa}
\usepackage{xfrac}
\usepackage[unicode,hidelinks,bookmarks=false]{hyperref}
\newcommand{\ie}{i.e.\ }
\newcommand{\cf}{cf.\ }
\newcommand{\resp}{respectively\ }
\newcommand{\Subsection}{\subsection}
\providecommand{\nobreakdash}{\nobreak\hbox{-}}
\setlength{\emergencystretch}{2em}
\multlinegap0pt
\usepackage{mathtools}
\usepackage{amssymb}
\usepackage{tikz}
\newtheorem{lemma}{Lemma}
\newtheorem{definition}{Definition}
\newtheorem{proposition}{Proposition}
\DeclareMathOperator{\LDmap}{\scalebox{0.8}{\(\xrightarrow{\mathrm{LD}}\)}} 
\newcommand{\N}{\mathbb N}
\newcommand{\R}{\mathbb R}
\newcommand{\UC}{\(\mathrm{(UC)}\)}
\DeclareMathOperator{\cK}{\mathcal{K}}
\DeclareMathOperator{\cP}{\mathcal{P}} 
\DeclareMathOperator{\cQ}{\mathcal{Q}} 
\DeclareMathOperator{\cpt}{\scalebox{0.8}{\(\xsubset{\mathrm{c}}\)}} 
\renewcommand{\epsilon}{\varepsilon}
\newcommand{\lang}{\mathscr{L}}
\newcommand{\sub}{\sigma}
\newcommand{\xsubset}[1]{\overset{#1}{\subset}}
\title{\bf Recognisability for generalised hierarchical pattern spaces of finite local complexity}
\author{James J. Walton}
\date{Source: arXiv:2509.21001v2 (CC BY 4.0)}
\begin{document}
\maketitle

\noindent\emph{Source and licence.} \bf Recognisability for generalised hierarchical pattern spaces of finite local complexity. Version arXiv:2509.21001v2 (CC BY 4.0), distributed by arXiv under the Creative Commons Attribution 4.0 International licence, http://creativecommons.org/licenses/by/4.0/, as recorded in the arXiv licence metadata for this exact version and shown on the abstract page https://arxiv.org/abs/2509.21001v2. Full licence notice: \texttt{licenses/arxiv-2509.21001-CC-BY-4.0.txt}.\par\smallskip

\noindent\textit{Editorial scope.} Counted unit: the article's proposition prop:L-sub=>Pf (source0.tex lines 1653--1655). The article's complete original proof of it is delivered with it (source0.tex lines 1657--1716, one \texttt{proof} environment of 60 lines). No mathematics is written, repaired or re-derived: the statement and the proof are the article's own lines, in the article's own notation, and no retained line is substituted or re-flowed.  Retained uncounted as the source-local context the proof needs: lines 340--351; lines 565--567; lines 732--734; lines 1294--1307; lines 1498--1502; lines 1622--1631.  Nothing was omitted from any retained span, so the package carries no omission record.  No retained line contains a citation, so the wrapper carries no bibliography and no bibliography entry.  No retained line contains a figure environment, an image-inclusion command or drawing code, so no image is delivered and the package declares no asset dependency.  Editorial formatting only, in addition: an AMS \texttt{amsart} preamble that restates the article's own declarations and macros used by the retained lines, namely the environments \texttt{lemma}, \texttt{definition}, \texttt{proposition} on separate counters, together with the standard packages those lines need.  The retained environments print in this document as Lemma 1, Definition 1, Proposition 1 in order of appearance, whereas the article prints the same items with the numbering of its own full document.  The complete native source of the article as distributed in the arXiv e-print for this exact version is bundled unchanged as \texttt{sources/185720/source0.tex}.
\begin{lemma}\label{lem:pattern properties}
For patterns \(\cP\), \(\cQ \in A^E\),
\begin{enumerate}
	\item if \(\cP[x,V] = \cQ[y,V]\) and \(U \subseteq V\) then \(\cP[x,U] = \cQ[y,U]\);
	\item if \(T(x) = L(x) + z\) is an affine automorphism then \(\cP[x,V] = \cQ[y,V]\) if and only if \((T\cP)[Tx,LV] = (T\cQ)[Ty,LV]\). For \(z \in E\), we have that
	\[
	\cP[x,V] = \cQ[y,V] \text{ if and only if } \cP[x+z,V-z] = \cQ[y+z,V-z];
	\]
	\item if \(\cP[x,V] = \cQ[y,V]\) and \(U+z \subseteq V\) then \(\cP[x+z,U] = \cQ[y+z,U]\);
	\item if \(V = \bigcup_{i \in \mathcal{I}} U_i\) then \(\cP[x,V] = \cQ[y,V]\) if and only if \(\cP[x,U_i] = \cQ[y,U_i]\) for all \(i \in \mathcal{I}\).
\end{enumerate}
\end{lemma}

\begin{lemma}\label{lem:close repeat => period}
Let \(U \cpt E\) with \(\mathbf{0} \in U\). Suppose there exists \(x \in E\) and \(V \cpt E\) satisfying the following: every \(U\)-patch occurs with centre within \(V\) of a point \(x\). That is, suppose that for all \(p \in \lang_U(\cP)\) there exists some \(v \in V\) with \(\cP[x+v,U] = p\). Then, if \(\cP[x,V] = \cP[x+u,V]\) for some \(u \in U\), we have \(u \in \cK_{\cP}\).
\end{lemma}

\begin{definition}\label{def:FLC pattern space}
A pattern space \(\Omega\) is of \textbf{finite local complexity} (or is \textbf{FLC}) if, for each \(U \cpt E\), there are finitely many patterns \(\cP_1\), \(\cP_2\), \ldots, \(\cP_n\) (\(n\) possibly depending on \(U\)) and some \(K(U) \cpt E\) for which, for all \(p \in \lang_U(\Omega)\), there exists some \(i \in \{1,\ldots,n\}\) and \(v \in K(U)\) with \(p = \cP_i[v,U]\).
\end{definition}

\begin{lemma}\label{lem:bigger agreement after substitution}
Let \(\Omega\), \(\Omega'\) be pattern spaces and \(S \colon L\Omega' \LDmap \Omega\). As usual, define \(\sub \coloneqq S \circ L \colon \Omega' \to \Omega\). Then, for all \(\cP\), \(\cQ \in \Omega'\), we have that
\[
\mathrm{if} \ \cP[x,U] = \cQ[y,U] \ \mathrm{ then } \ (\sub \cP)[Lx,(LU)_{-c}] = \newline (\sub \cQ)[Ly,(LU)_{-c}] .
\]
In particular, when \(L\) is expansive, for all \(r \geq c/\lambda\) we have
\[
\mathrm{if} \ \cP[x,r] = \cQ[y,r] \ \mathrm{ then } \ (\sub \cP)[Lx,\lambda r - c] = (\sub \cQ)[Ly,\lambda r - c] 
\]
and, for all \(k \in \N_0\),
\[
\mathrm{if} \  \cP[x,V^k] = \cQ[y,V^k] \ \mathrm{ then } \ (\sub \cP)[Lx,V^{k+1}] = (\sub \cQ)[Ly,V^{k+1}] .
\]
\end{lemma}

\begin{proposition}\label{prop:full patch at origin}
Let \(\Omega\) be compact Hausdorff and expansive \(L\)-sub. There exists some \(N \in \N\) so that \(\sub^N\) fixes LI-classes and \(Y \cpt E\) satisfying the following. Let \(\cP \in \Omega\) be arbitrary and choose any hierarchy \((\cP_i)_{i=0}^\infty\) of \(\cP\). Then, for all \(i \in \N_0\), every \(V^0\)-patch of \(\cP\) appears in \(\cP_{Ni}\) with centre in \(Y\). That is, for all \(i \in \N_0\) and \(p \in \lang_0(\cP)\), there exists some \(y \in Y\) with \(\cP_{Ni}[y,V^0] = p\).

Thus, there exists \(\ell \in \N\) satisfying the following. For all \(\cP \in \Omega\) and \(i \in \N_0\), every \(V^i\)-patch of \(\cP\) appears with centre in \(V^{i+\ell}\) in \(\cP\). That is, for all \(p \in \lang_i(\cP)\), there exists some \(y \in V^{i+\ell}\) with \(p = \cP[y,V^i]\).
\end{proposition}

\begin{lemma}\label{lem:local recognisability}
Given \(\cP \in \Omega\), \(\cP' \in \sub^{-1}(\cP)\) and \(f \colon \R_{\geq 0} \to \R_{\geq 0}\), define the following property:
\begin{itemize}
	\item[(\(\mathrm{P}_f\))] for all \(x \in E\) and \(R \geq 0\),
	\[
	\text{if} \ \cP[\mathbf{0},f(R)] = \cP[x,f(R)] \ \text{then there exists } g \in \cK_{\cP} \ \text{with} \ \cP'[\mathbf{0},R] = \cP'[L^{-1}(x+g),R] . 
	\]
\end{itemize}
If, for all \(\cP \in \Omega\) and \(\cP' \in \sub^{-1}(\cP)\), there exists some \(f\) for which \(\mathrm{P}_f\) holds for this pair, then \(\sub\) satisfies \UC.
\end{lemma}

\begin{proposition}\label{prop:L-sub=>Pf}
There exists some \(f \colon \R_{\geq 0} \to \R_{\geq 0}\) for which Property \(\mathrm{P}_f\) of Lemma \ref{lem:local recognisability} holds for all \(\cP \in \Omega\) and \(\cP' \in \sub^{-1}(\cP)\).
\end{proposition}

\begin{proof}
Let \(\cP \in \Omega\) and \(\cP' \in \sub^{-1}(\cP)\) be arbitrary and define \(\cK \coloneqq \cK_{\cP}\). Choose a corresponding hierarchy \((\cP_i)_{i=0}^\infty\), that is, a sequence of patterns \(\cP_i \in \Omega\) with \(\sub(\cP_{i+1}) = \cP_i\), \(\cP_0 = \cP\) and \(\cP_1 = \cP'\). Throughout this proof, to simplify notation, for \(j \in \N_0\) we denote \(\cP_i[x,j] \coloneqq \cP_i[x,V^j]\) (instead of our usual notation \(\cP_i[x,j] = \cP_i[x,B_j]\)). We claim that, for sufficiently large \(k \in \N_0\), if \(\cP[\mathbf{0},k+2] = \cP[x,k+2]\) for some \(x \in E\) then \(\cP'[\mathbf{0},k] = \cP'[x^{-1} + L^{-1} h,k]\) for some \(h \in \cK\). In this case, clearly \(\mathrm{P}_f\) holds (and even independently of \(\cP\) and \(\cP'\)) for a certain function \(f\) (whose explicit description is not important and is neater in terms of \(V^k\)-patches rather than \(r\)-patches).

Let \(\ell \in \N\) be as given from Proposition \ref{prop:full patch at origin}. Let \(\epsilon > 0\) with \(\epsilon \leq (2(\overline{\lambda})^\ell)^{-1}\), where we take \(\overline{\lambda} > 0\) such that \(LB \subseteq \overline{\lambda} B\) for \(B \coloneqq B_1\). We may also take \(\epsilon > 0\) sufficiently small that \(V^k + L^k(\epsilon B) \subseteq V^{k+1}\) for all \(k \in \N_0\), since
\begin{align*}
V^{k+1} &= L^{k+1}(B) + \kappa B \supseteq \lambda L^k(B) + \kappa B = (L^k(B) + (\lambda - 1)L^k(B)) + \kappa B\\
        &= V^k + L^k((\lambda - 1)B) \supseteq V^k + L^k(\epsilon B) ,
\end{align*}
provided \(0 < \epsilon \leq \lambda - 1\). In the second equality, we use that for all \(\alpha\), \(\beta > 0\),
\[
\alpha L^k(B) + \beta L^k(B) = L^k(\alpha B + \beta B) = L^k((\alpha + \beta)B) = (\alpha + \beta)L^k(B) .
\]
By expansivity, we also have \(V^k + L^m(\epsilon B) \subseteq V^{k+1}\) for other \(m \leq k\).

From compactness of \(\Omega\), we may apply FLC (Definition \ref{def:FLC pattern space}) to patches with shape \(U = V^3\), yielding patterns \(\cQ_1\), \ldots, \(\cQ_n \in \Omega\) and \(K \cpt E\) for which, for all \(\cQ \in \Omega\) and \(x \in E\), there exists some \(i \in \{1,\ldots,n\}\) and \(x \in K\) with \(\cQ[x,3] = \cQ_i[y,3]\). Let \(K' = K \times \{1,\ldots,n\}\) denote the disjoint union of \(n\) copies of \(K\). By compactness of \(K' \times K'\), there is some \(M \in \N\) for which, for any \(M\) pairs \(((a_i,m_i),(b_i,n_i)) \in K' \times K'\), we may find \(i \neq j\) with \(m_i = m_j\), \(n_i = n_j\), \(\|a_i - a_j\| \leq \epsilon\) and \(\|b_i - b_j\| \leq \epsilon\).

For a contradiction, suppose that there are infinitely many \(k \in \N_0\) and \(x_k \in E\) with \(\cP[\mathbf{0},k+2] = \cP[x_k,k+2]\) yet \(\cP'[\mathbf{0},k] \neq \cP'[L^{-1}(x_k + g),k]\) for all \(g \in \cK\). For each such \(k\), consider the patches \(\cP_k[\mathbf{0},3]\) and \(\cP_k[L^{-k}x_k,3]\). From FLC, we may write \(\cP_k[\mathbf{0},3] = \cQ_{m_k}[a_k,3]\) and \(\cP_k[L^{-k} x_k,3] = \cQ_{n_k}[b_k,3]\) for \(a_k\), \(b_k \in K\) and \(m_k\), \(n_k \in \{1,\ldots,n\}\). By considering \(M\) different values of \(k \geq \ell\), we obtain corresponding pairs \((a_k,m_k) \in K'\) and \((b_k,n_k) \in K'\). By our choice of \(M\), there exist indices \(j > i \geq \ell\) for which \(m_i = m_j = m\), \(n_i = n_j = n\), \(\|a_i - a_j\| \leq \epsilon\) and \(\|b_i - b_j\| \leq \epsilon\). In summary,
\[
\cP_i[\mathbf{0},3] = \cQ_m[a_i,3]\ , \ \cP_i[L^{-i} x_i,3] = \cQ_n[b_i,3]\ , \ \cP_j[\mathbf{0},3] = \cQ_m[a_j,3]\ , \ \cP_j[L^{-j} x_j,3] = \cQ_n[b_j,3].
\]
By shifting the third equation above by \(u \coloneqq a_i - a_j \in \epsilon B\) and the fourth by \(v \coloneqq b_i - b_j \in \epsilon B\) (using Lemma \ref{lem:pattern properties} (3), with \(V^2 + \epsilon B \subseteq V^3\)), and combining with the first and second equations, respectively (each restricted to \(V^2\)), we obtain
\begin{equation}\label{eq:UC1}
\cP_i[\mathbf{0},2] = \cP_j[u,2]\ , \ \cP_i[L^{-i}x,2] = \cP_j[L^{-j}y + v,2] ,
\end{equation}
where we denote \(x \coloneqq x_i\) and \(y \coloneqq x_j\). We recall that, by definition of \(x_i\) and \(x_j\),
\begin{equation}\label{eq:UC2}
\cP[\mathbf{0},i+2] = \cP[x,i+2]\ , \ \cP[\mathbf{0},j+2] = \cP[y,j+2]
\end{equation}
but, for all \(g \in \cK\),
\begin{equation}\label{eq:UC3}
\cP'[\mathbf{0},i] \neq \cP'[L^{-1}(x+g),i]\ , \ \cP'[\mathbf{0},j] \neq \cP'[L^{-1}(y+g),j] .
\end{equation}
We will construct some \(h \in \cK\) contradicting the right-hand non-equality for \(g=-h \in \cK\).

Begin by substituting Equation \ref{eq:UC1} \(i\) times, by Lemma \ref{lem:bigger agreement after substitution}, and denoting \(\eta \coloneqq j-i \in \N\), to obtain
\begin{equation}\label{eq:UC4}
\cP[\mathbf{0},i+2] = \cP_\eta[L^i u,i+2]\ , \ \cP[x,i+2] = \cP_\eta[L^{-\eta} y + L^i v,i+2] .
\end{equation}
Combining this with Equation \ref{eq:UC2} gives
\begin{equation}\label{eq:UC5}
\cP_\eta[L^i u,i+2] = \cP[\mathbf{0},i+2] = \cP[x,i+2] = \cP_\eta[L^{-\eta} y + L^i v,i+2] .
\end{equation}
Substituting this a further \(\eta\) times gives
\[
\cP[L^j u,j+2] = \cP[y + L^j v,j+2] .
\]
Recall that \(u \in \epsilon B\) and that, for all \(j \in \N_0\), we have \(V^{j+2} \supseteq V^{j+1} + L^j(\epsilon B)\). Thus, we may shift the above equation by \(L^j v \in L^j(\epsilon B)\) to obtain
\begin{equation}
\cP[L^j(u-v),j+1] = \cP[y,j+1] = \cP[\mathbf{0},j+1] ,
\end{equation}
where the latter is given by restricting the right-hand side of Equation \ref{eq:UC2} to \(V^{j+1}\). The above shows that \(h \coloneqq L^j(u-v)\) is a `short' return vector to the relatively large patch of shape \(V^j\) (in fact, even \(V^{j+1}\)). More precisely, we have \(h \in V^{j-\ell}\) since, by assumption, \(h = L^j(v - u) \in L^j(2\epsilon B)\), so
\[
h \in L^{j-\ell}L^\ell(2\epsilon B) \subseteq L^{j-\ell}((\overline{\lambda})^\ell 2\epsilon B) \subseteq L^{j-\ell} B \subseteq L^{j-\ell} B + \kappa B = V^{j-\ell} ,
\]
where we recall, for the second inclusion, that we defined \(\epsilon \leq (2(\overline{\lambda})^\ell)^{-1}\). In the first, we use that \(L(B) \subseteq \overline{\lambda}B\), hence \(L(\mu B) \subseteq \overline{\lambda}(\mu B)\) for any \(\mu \geq 0\) (here, \(\mu = 2\epsilon\)) and thus, by induction, \(L^\ell(\mu B) \subseteq (\overline{\lambda})^\ell \mu B\).

Summarising, we have \(h = L^j(u-v) \in V^{j-\ell}\) for which \(\cP[h,j+1] = \cP[\mathbf{0},j+1]\). Recall that \(\ell\) was chosen from Proposition \ref{prop:full patch at origin} so that a centre of every \(V^n\)-patch occurs in \(V^{n+\ell}\), that is, given \(p \in \lang_n(\cP)\), we have that \(p = \cP[x,V^n]\) for some \(x \in V^{n+\ell}\). Taking \(n = j-\ell \in \N_0\) here, since \(h \in V^{j-\ell}\), it follows from Lemma \ref{lem:close repeat => period} (in short: a \(U\)-short return vectors of a patch containing centres of all \(U\)-patches must be a period) that \(h \in \cK\).

Thus, it remains to show that \(\cP'[\mathbf{0},j] = \cP'[L^{-1}(y-h),j]\), contradicting the assumed non-equality of Equation \ref{eq:UC3}. But this may be done similarly to earlier, just shifting down \(j-1\) rather than \(j\) levels. Explicitly, from Equation \ref{eq:UC5}, substitute a further \(\eta-1\) times to obtain \(\cP_1[L^{j-1} u,j+1] = \cP_1[L^{-1} y + L^{j-1} v,j+1]\). Just as before, due to our choice of small \(\epsilon\), we have \(V^{j+1} \supseteq V^j + L^{j-1}(\epsilon B)\), so we may shift this equality of patches by \(-L^{j-1} u \in L^{j-1}(\epsilon B)\) to obtain \(\cP_1[\mathbf{0},j] = \cP_1[L^{-1} y + L^{j-1}(v-u),j]\). Since \(L^{j-1}(v-u) = L^{-1}(L^j(v-u)) = L^{-1}(-h)\), the result follows.
\end{proof}

\end{document}
